
\section{Updated Proposal}

We take a graph, upon which there are players who wish to identify if the
graph possesses some property. We wish to use this model to find the minimum
communication required for each player to identify if this property exists.
Each Player will receive some portion of the graph as their input--this can be
either vertices, edges, or both. We mostly study reachability of vertices, in
which the only paths considered valid are paths which use only the portions of
the graph given to the players as inputs.

For motivation, please see Section \ref{sec:introduction} of the report.

\subsection{Project Goals}

Sections \ref{sec:two-party}, \ref{sec:iter-two-party}, \ref{sec:multi-party}, and \ref{sec:iter-multi-party} formally define our problem statement and the full spectrum of problems we are considering for our project.

Clearly, section \ref{sec:two-party} is done. Much of the results for the remaining sections are also complete, and it only remains to formalize and type them up in \LaTeX.

\subsection{Actionables}

Therefore, this project can be summarized into three specific, actionables:

\begin{enumerate}
    \item Finish writing this paper to outline the formal, technical communication complexity bounds for a given multi-agent traversal problem with certain constraints.
    \item Finish porting the report to our website, whose link is available on the cover page of this report.
    \item Give a poster presentation of the key results of the paper, with the intention of making the results applicable to a wider audience.
\end{enumerate}

Therefore, our work will remain theoretically grounded while remaining applicable for the wider computing community.

\subsection{Deliverables}

We have three main deliverables:

\begin{enumerate}
    \item A paper presenting actionable item 1 and 2.
    \item A research poster of our key results, for actionable 3.
    \item A website containing the paper and poster information, for actionable 3.
\end{enumerate}

To be more specific, the paper will take the form of a research paper, where we define the problems we are working on, prove bounds on these problems, and show how they model situations in the field of multi-agent graph traversal. The paper is our primary project output.

Our detailed weekly progress is shown in the previous appendix section. In the remaining time, we plan for the following actionables for each person:

\small
\setlength{\tabcolsep}{4pt}
\begin{longtable}{@{}>{\raggedright\arraybackslash}p{0.06\textwidth}>{\raggedright\arraybackslash}p{0.22\textwidth}>{\raggedright\arraybackslash}p{0.22\textwidth}>{\raggedright\arraybackslash}p{0.22\textwidth}>{\raggedright\arraybackslash}p{0.22\textwidth}@{}}
\caption{Weekly plan }
\label{tab:weekly-contributions}\\
\toprule
Week & Ryan Batubara & Ivy Hawks & Darren Ho & Ciro Zhang \\
\midrule
\endfirsthead
\toprule
Week & Ryan Batubara & Ivy Hawks & Darren Ho & Ciro Zhang \\
\midrule
\endhead
\bottomrule
\endfoot
\bottomrule
\endlastfoot
7 
& Finish the other problems section, with a focus on figures.
& Finish the iterated multi-party section, with a focus on figures.
& Make significant progress on the multi party section, with a focus on figures.
& Make significant progress on the iterated multi-party section, with a focus on figures.
\\
8 
& Finish the multi-party section.
& Finish the iterated multi-party section.
& Finish the website port.
& Move figures to first draft of poster.
\\
9
& Polish paper and proofread website.
& Polish website and proofread poster.
& Polish poster and proofread paper.
& Polish poster and proofread paper.
\end{longtable}
\normalsize

Indeed, given our weekly contributions thus far we are on track to finishing the project with a week dedicated solely to polishing.
